% Header: Here are all packages used and some additional definitions
%%%%%%%%%%%%%%%%%%%%%%%%%%%%%%%%%%%%%%%%%%%%%%%%%%%%%%%%%%%%%%%%%%%

\documentclass[11pt,a4paper]{scrartcl}
\usepackage[margin=2.5cm]{geometry}
\usepackage[onehalfspacing]{setspace}
\usepackage{graphicx} % zum Einbinden von Graphiken
\usepackage{enumitem}
\usepackage{tcolorbox}
\tcbuselibrary{breakable}
\usepackage[breaklinks=true,colorlinks=true,linkcolor=blue,urlcolor=blue,citecolor=blue]{hyperref} % f. Referenzen
\usepackage{amsmath,amsthm,amssymb} % Mathematik Umgebung 
\usepackage{icomma} % Intelligentes Komma, das den richtigen Abstand zwischen Dezimalzahlen als auch in Formeln wählt.
\usepackage[ngerman]{babel} % Deutsche Bezeichnungen bei Inhaltsangabe etc
\usepackage[T1]{fontenc}    % andere Schriftsatzkodierung für richtige Silbentrennung bei Umlauten
\usepackage[locale = DE,space-before-unit=true,per-mode = symbol]{siunitx} % Bessere Einheiten
\usepackage{booktabs,multirow} % Pakete zur Erstellung von Tabellen
\usepackage{placeins} % Definiert den Befehl “\FloatBarrier”, der die Ausgabe der davor eingebundenen Bilder erzwingt, befor der Text weiter geht. (Mit vorsicht zu verwenden)
\usepackage[natbib,abbreviate=true,doi=false,style=numeric-comp,giveninits=true,sorting=none]{biblatex} % Modernes Paket zur Erzeugung von Bibliografien (benötigt biber!)
\usepackage{csquotes} % Fortgeschrittene Funktionen für Zitate, für die deutsche Form der Anführungszeichen bei Referenzen
%\addbibresource{MyBibliography.bib} % Ort der .bib Datei, die die Datenbank für Literatur/Referenzen enthält.

\graphicspath{{Bilder/} {images/} {chapters/}}

\DeclareSIUnit{\dBm}{dBm}
\DeclareSIUnit[per-mode=reciprocal]\WN{\per\centi\meter}

%%%%%%%%%%%%%%%%%%%%%%%%%%%%%%%%%%%%%%%%%%%%%%%%%%%%%%%%%%%%%%%%%%%
\begin{document}
%
\title{Experimental Physics - Musterlösung Übungsblatt 2 Aufgabe 1}
\maketitle
\vfill
\thispagestyle{empty}
%
%
\thispagestyle{empty}
\pagenumbering{arabic} 
%
%
\paragraph{Aufgabe 1: Formfaktor - Herleitung}
In der Vorlesung wurde der Formfaktor für die Ladungsverteilung $f(x)$ eingeführt, der als 
\begin{align*}
    F(\vec{q}) = \int e^{i \, \vec{q} \cdot \vec{x}/\hbar} f(\vec{x})\, d^3x
\end{align*}
definiert ist.
\begin{enumerate}[label=(\alph*)]
    \item Im Folgenden werden nur kugelsymmetrische Ladungsverteilungen betrachtet. Zeigen Sie, daß in diesem Fall der Formfaktor zu
    \begin{align*}
        F(q) = 4\pi\int\limits_0^\infty \frac{\sin{q \, r/\hbar}}{q \, r / \hbar} f(r) \, r^2 dr
    \end{align*}
    umgeformt werden kann (F hängt also nur vom Betrag von $\vec{q}$ ab).\\
    \underline{Hinweis:} Wählen sie das Koordinatensystem so, daß $q$ in Richtung der z-Achste weist.
    
    \begin{tcolorbox}[title=Berechnung, breakable]
        Ausgehend von der Definition des Formfaktors
        \begin{align*}
            F(\vec{q}) = \int d^3x \, e^{i \, \vec{q} \cdot \vec{x}/\hbar} f(\vec{x})
        \end{align*}
        verwende man Kugelkoordinaten
        \begin{align*}
            \vec{x} = \begin{bmatrix}
                x\\y\\z
            \end{bmatrix}=\begin{bmatrix}
                r \sin{\theta} \cos{\phi}\\
                r \sin{\theta} \sin{\phi}\\
                r \cos{\theta}
            \end{bmatrix},              
        \end{align*}
        den hingewiesenen Ausdruck
        \begin{align*}
            \vec{q} = \begin{bmatrix}
                0\\0\\q
            \end{bmatrix} \implies \vec{q} \cdot \vec{x} = q \, r \cos{\theta}
        \end{align*}
        und den Volumenselement ($J$ die Jakobi-Matrix; Ermittlung in den Nebenrechnungen am Ende)
        \begin{align*}
            d^3x = dx\,dy\,dz = |\det{J}|\,dr\,d\theta\,d\phi = r^2 \sin{\theta}\, dr\, d\theta \, d\phi
        \end{align*}

        um auf folgenden Ausdruck zu kommen:

        \begin{align*}
            \int\limits_0^\infty dr \int\limits_0^{\pi}d\theta \int\limits_0^{2\pi} \, d\phi \, e^{i \, q \, r/\hbar\cos{\theta}} f(r) \, r^2 \sin{\theta}.
        \end{align*}
        Das Integral über $\phi$ gibt $2\pi$ und mit folgender Substitution
        \begin{align*}
            u &:= \cos{\theta},\, u(0) = 1, \, u(\pi) = -1\\
            \frac{du}{d\theta} &= -\sin{\theta},\, du = -\sin{\theta} \, d\theta
        \end{align*}
        kann man das Integral über $\theta$ lösen:
        \begin{align*}
            - 2 \pi \int\limits_0^\infty dr &\int\limits_{+1}^{-1}du \, e^{i \, q \, r/\hbar u} f(r) \, r^2=\\ 
            - 2 \pi \int\limits_0^\infty dr &\left[\frac{e^{i \, q \, r/\hbar u}}{i \, q \, r/\hbar}\right]_{+1}^{-1} f(r) \, r^2 =\\
            2 \pi \int\limits_0^\infty dr &\left(\frac{e^{+i \, q \, r/\hbar} - e^{-i \, q \, r/\hbar}}{i \, q \, r/\hbar}\right) f(r) \, r^2.
        \end{align*}
        Mit Verwendung des trigonometrischen Zusammenhangs mit der komplexen Exponentialfunktion (Nebenrechnung) ergibt sich der gesuchte Ausdruck:
        \begin{align*}
            F(q) = 4 \pi \int\limits_0^\infty dr &\frac{\sin(q \, r/\hbar)}{q \, r/\hbar} f(r) \, r^2.
        \end{align*}
    \end{tcolorbox}

\end{enumerate}
Nun soll ein Spezialfall betrachtet werden, auf den wir in späteren Augaben noch zurückkommen werden.
\begin{enumerate}[start=2, label=(\alph*)]
    \item Zeigen Sie, dass für den Formfaktor einer homogen geladenen Kugel:
    \begin{align*}
        F(q) = 3\frac{\sin{u} - u \cos{u}}{u^3}
    \end{align*}
    gilt und bestimmen Sie die Form von u. Für welche Teilchen gilt eine solche Ladungsverteilung in guter Näherung?\\
    \underline{Hinweis:} Überlegen sie dazu zuerst, wie in diesem Fall die Ladungsdichteverteilung aussieht, wenn Sie berücksichtigen, dass diese über
    \begin{align*}
        \int f(x) d^3x = 1
    \end{align*}
    normiert ist.

    
    \begin{tcolorbox}[title=Berechnung, breakable]
        Zunächst wird versucht $f(r)$ zu finden. Die Ladungsdichte einer homogen geladene Kugel ist innerhalb der Kugel konstant und außerhalb null. 
        \begin{align*}
            f(r) =
            \begin{cases}
                C, & 0 \leq r \leq R, \\
                0, & r > R .
            \end{cases}
        \end{align*}
        Aus der Normierung kann die Konstante $C$ ermittelt werden ($4\pi$ aus Kugelkoordinaten):
        \begin{align*}
            1 = \int d^3x\, f(x) = 4\pi \int_0^{\infty} dr\, f(r)\, r^2 = 4\pi \int_0^R dr\, C\, r^2 = 4\pi C\left[\frac{r^3}{3}\right]_0^R = 4\pi C\frac{R^3}{3}\\
            C = \frac{3}{4\pi R^3}
        \end{align*}
        Damit hat man die Ladungsdichte
        \begin{align*}
            f(r) =
            \begin{cases}
                \frac{3}{4\pi R^3}, & 0 \leq r \leq R, \\
                0, & r > R .
            \end{cases}
        \end{align*}
        Dann setze man die Ladungsdichte in das vorher ermittelte Integral:
        \begin{align*}
            F(q) &= 4 \pi \int\limits_0^\infty dr \frac{\sin(q \, r/\hbar)}{q \, r/\hbar} f(r) \, r^2 \\
            &= 4 \pi \int\limits_0^R dr \frac{\sin(q \, r/\hbar)}{q \, r/\hbar} \frac{3}{4\pi R^3} \, r^2\\
            &= \frac{3}{R^3} \int\limits_0^R dr \frac{\sin(q \, r/\hbar)}{q \, r/\hbar}  \, r^2\\
        \end{align*}
        Wenden hier partielle Integration an
        \begin{align*}
            \frac{3}{R^3 q/\hbar} \int\limits_0^R dr \sin(q \, r/\hbar)  \, r &= \frac{3}{R^3 q /\hbar}  \left(\left[\frac{r(-\cos(q\,r/\hbar))}{q/\hbar}\right]_0^R - \int\limits_0^R dr \frac{(-\cos(q\,r/\hbar))}{q/\hbar}\right)\\
            &= \frac{3}{R^3 q /\hbar}  \left[-\frac{r\cos(q\,r/\hbar)}{q/\hbar} + \frac{\sin(q\,r/\hbar)}{(q/\hbar)^2}\right]_0^R\\
            &= \frac{3}{R^3 q /\hbar}  \left(-\frac{R\cos(q\,R/\hbar)}{q/\hbar} + \frac{\sin(q\,R/\hbar)}{(q/\hbar)^2}\right)
        \end{align*}
        Etwas Umstellen und man bekommt die gesuchte form mit $u = q\,R/\hbar$
        \begin{align*}
            &= 3 \left(\frac{\sin(q\,R/\hbar)}{R^3(q/\hbar)^3} -\frac{R\cos(q\,R/\hbar)}{R^3 (q/\hbar)^2}\right)\\
            &= 3 \frac{\sin(q\,R/\hbar) - q\,R/\hbar\cos(q\,R/\hbar)}{(q\,R/\hbar)^3} \\
            F(q) &= 3 \frac{\sin{u} - u \cos{u}}{u^3}
        \end{align*}
    \end{tcolorbox}

    \begin{tcolorbox}[title=Antwort auf Frage, breakable]
        
    \end{tcolorbox}

    \item Welcher Wert ergibt sich für den Formfaktor aus (a) für $q=0$?

    \begin{tcolorbox}[title=Berechnung, breakable]
    \end{tcolorbox}
\end{enumerate}

\textbf{Nebenrechnungen}
\begin{tcolorbox}[title=Nebenrechnung, breakable]
\end{tcolorbox}

\end{document}
